\section{Бёрч–Свиннертон-Дайер и соседи за горизонтом}\label{sec:bsd}

Восьмая маска, гипотеза Бёрча–Свиннертон-Дайера~\cite{clay-bsd}, входит в семью иначе,
чем семь предыдущих. Там движок \emph{сторожит}: отклонение от нормы — замаскированный вечный
двигатель, и запрет его убивает. У BSD движок не сторожит — он \emph{работает методом}.
Алгебраический ранг эллиптической кривой доказан конечным ровно бесконечным спуском Ферма, тем самым
спуском, что лежит в основании всей программы. Здесь он не запрет, а инструмент. Lean-источник —
\srcpath{Engine/BSDFront.lean}; ядро заземлено на реальные объекты \code{mathlib}
(\code{WeierstrassCurve.Affine.Point}, \code{WeierstrassCurve.LSeries}, класс \code{Northcott},
\code{AddCommGroup.fg\_of\_descent}). Это \emph{не} доказательство BSD.

\subsection{Спуск — это и есть движок}

Теорема Морделла–Вейля говорит: группа рациональных точек $E(\mathbb{Q})$ конечно порождена,
$E(\mathbb{Q})\cong\mathbb{Z}^r\oplus T$. Доказывают её спуском: каноническая высота Нерона–Тейта —
положительно определённая квадратичная форма, а по свойству Нортокотта точек ограниченной высоты
конечно много, значит спуск по высоте обязан оборваться. Это ровно наш запрет: бесконечная строго
убывающая цепь высот — вечный двигатель, а его нет.

\begin{theorem}[\code{bsd\_descent\_has\_no\_engine}]\label{thm:bsd_descent_has_no_engine}
\green{} На реальной группе точек $W.\code{toAffine.Point}$ эллиптической кривой высотный спуск не
имеет вечного двигателя: для любой высотной модели спуска $M$ не существует
$\code{PerpetualEngine}\,M.\code{descentStep}$, то есть строго убывающая по высоте
$h\colon W.\code{toAffine.Point}\to\mathbb{N}$ бесконечная цепь точек невозможна.
\end{theorem}

\emph{Идея доказательства.} Высота — это $\mathbb{N}$-ранг, а ранг переносит запрет двигателя из
универсальной формы через \code{UniversalEngine.no\_perpetual\_engine\_of\_natRank}
(\cref{sec:engine}). Носитель — настоящий \code{WeierstrassCurve.Affine.Point} с его групповым
законом, а механизм — тот же, что у \code{Northcott} и \code{AddCommGroup.fg\_of\_descent}:
фундированность высоты обрывает спуск. Модель обитаема (не вакуумна), с конкретным свидетелем над
$\mathbb{Q}$.

Мы честны о границе дозволенного. Каноническую высоту Нерона–Тейта \code{mathlib} ещё не содержит
(есть Вейль-высота и приближённый закон параллелограмма), поэтому высота у нас \emph{именована}. Но
носитель и групповой закон реальны, и запрет спуска выведен строго: тень лежит на настоящей кривой,
а не на макете.

\subsection{Паритет — тот же rank-parity узел}

Гипотеза о паритете утверждает $(-1)^{\operatorname{rank}}=w(E)$, где $w(E)$ — корневое число, знак
функционального уравнения. А $(-1)^{\operatorname{rank}}$ — в точности тот rank-parity инвариант,
что стоит за Риманом: функция Лиувилля $\lambda(n)=(-1)^{\Omega(n)}$.

\begin{theorem}[\code{bsd\_parity\_is\_rankParity}]\label{thm:bsd_parity_is_rankParity}
\green{} Чётность ранга совпадает с лиувиллевой: для любых $r,n\in\mathbb{N}$ с $n\neq 0$ и
$\Omega(n)=r$ верно $(-1)^r=\lambda(n)$ в $\mathbb{Z}$.
\end{theorem}

\emph{Идея доказательства.} Прямой мост к
\code{RiemannLiouville.liouville\_eq\_neg\_one\_pow\_rank} (\cref{sec:riemann}) — то же знаковое
правило, тот же флип при отнятии простого множителя. Паритет BSD — не новый зверь, а наш узел,
надетый на ранг эллиптической кривой.

Не декретировать ли этот паритет первопричиной, как мы декретировали закон Римана? Мы проверили — и
\emph{нет}. Настоящее корневое число $w(E)$ — глубокий аналитический инвариант, которого в
\code{mathlib} нет. Над доступной заглушкой (\code{RootNumber}\,$w:=w=\pm1$, свободное значение)
закон $(-1)^r=w$ удовлетворяется выбором $w$ (\code{bsd\_parityLaw\_satisfiable}), а универсальная
форма рефутируема (\code{bsd\_parityLaw\_not\_universal}). Декрет паритета был бы поэтому
\emph{вакуумен} — как у P/NP, Янг–Миллса и Ходжа. Мы не трогаем аксиому и оставляем паритет честно
открытым: связь с узлом концептуальна, а не декрет.

\subsection{Аналитический мост — честно открытый}

Сердце BSD — равенство алгебраического и аналитического рангов,
$\operatorname{rank}E(\mathbb{Q})=\operatorname{ord}_{s=1}L(E,s)$. \code{mathlib} уже определяет
$L$-функцию кривой (\code{WeierstrassCurve.LSeries}, через Эйлерово произведение), и мы на неё честно
ссылаемся. Но её аналитические свойства — продолжение к $s=1$, порядок нуля, функциональное
уравнение — не доказаны нигде.

\begin{definition}[\code{WeakBSD}]\label{def:WeakBSD}
\red{} Открытый вход: для эллиптической кривой $W/K$ и натуральных $\text{algRank},\text{aRank}$
предикат $\operatorname{WeakBSD}(W,\text{algRank},\text{aRank})$ есть
$\operatorname{AnalyticRank}(W,\text{aRank})\wedge\text{algRank}=\text{aRank}$. Мы его не
доказываем: это именованный вход, а не теорема.
\end{definition}

Люди прошли лишь края: при аналитическом ранге $\le 1$ (Гросс–Загир и Колывагин, через точки Хегнера
и эйлеровы системы) BSD доказана и Ша конечна; в среднем — свыше $66\%$ кривых
(Бхаргава–Скиннер–Чжан). При ранге $\ge 2$ открыто всё. Наш движок аналитику не закрывает: его роль —
спуск, а не мост.

У BSD нет собственной эпистемической оси, и это зафиксировано машинно
(\srcpath{Engine/BSDEpistemic.lean}): предложенная пара ground/beyondOwnHorizon вырождается, её вторая
нога — свободное поле убывания высоты, так что связка совпадает с уже доказанным запретом спуска
(\code{bsdGround\_coincides\_with\_engine}). Ворота \code{AnalyticRank} свободно выполнимы
(\code{analyticRank\_gate\_satisfiable}), а \code{WeakBSD} сводится к голому
$\text{algRank}=\text{aRank}$ (\code{weakBSD\_reduces\_to\_bare\_equality}) — именно там остаётся
честная \red{}-открытость. Содержательная замена (\srcpath{Engine/BSDAnalyticAnchor.lean}) делает
аналитический ранг настоящим понятием, $\code{analyticRankOf}\,L:=\code{analyticOrderAt}\,
L.\code{Lambda}\,1$, но лишь \emph{когда подана сама данность}: красный вход
\code{ContinuedLFunction} несёт $\Lambda$, совпадающую с \code{WeierstrassCurve.LSeries} на области
абсолютной суммируемости (\code{agreement}) и аналитичную в $s=1$ (\code{analyticAtOne}) — то самое
продолжение, которого в \code{mathlib} нет (даже с границей Хассе точка $1$ лежит вне полуплоскости
сходимости, $1<3/2$). На этих данных содержательные ворота \code{WeakBSDContentful} пришпиливают ранг
(\code{contentfulGate\_pins\_rank}), ценою данных, а не декрета.

Формализаций BSD, Морделла–Вейля, канонической высоты, аналитики $L$-функций эллиптических кривых нет
нигде по нашим разысканиям. Даже структурное engine/rank-parity прочтение, заземлённое на реальные
объекты \code{mathlib}, — первое в своём роде, но это \emph{не} решение: зелено лишь спуск и
паритет-узел, аналитический мост честно \red{}, декрета не добавлено.

\subsection{Пять открытых соседей}

Запрет вечного двигателя — спуск, высота, паритет ранга — живёт и в соседних открытых задачах. Мы
даём тень пяти из них одним приёмом: \emph{реальный якорь \code{mathlib} плюс честный гейт}. Якорь —
доказанная теорема (полиномиальный аналог, конечность, структура), которая и есть прочтение через
движок; гейт — сама открытая гипотеза, названная и не доказанная. Lean-источники:
\srcpath{Engine/ChowlaFront.lean}, \code{AbcFront.lean}, \code{BealFront.lean}, \code{LehmerFront.lean},
\code{LandauFront.lean}. Ни одна не решена; новизна формализационная.

\paragraph{Чоула и Сарнак.}
Функция Лиувилля $\lambda(n)=(-1)^{\Omega(n)}$ — наш rank-parity инвариант. Гипотеза Чоулы говорит:
паритет $\Omega$ не коррелирует по сдвигам, $\sum_{n\le x}\lambda(n)\lambda(n+h)=o(x)$ — тот же
запрет коллапса паритета в одну чётность, что охраняет и близнецов, и Римана.

\begin{theorem}[\code{chowlaCorrelation\_zero\_eq\_card}]\label{thm:chowlaCorrelation_zero_eq_card}
\green{} Диагональ $h=0$ даёт совершенную само-корреляцию:
$\sum_{n\le x}\lambda(n)^2=x$.
\end{theorem}

\emph{Идея доказательства.} $\lambda(n)^2=1$ для $n\ge 1$ (значения $\pm1$), так что диагональ
вырождается в подсчёт. Это резкий контраст тому, что Чоула запрещает при $h>0$, где сумма обязана
гаситься; флип знака при умножении на простое (\code{chowla\_parity\_flip}, реюз
\code{RiemannLiouville}) — тот же ход спуска. Гейты \code{ChowlaConjecture} и \code{SarnakConjecture}
(\red{}) — корреляции $o(x)$ и ортогональность Мёбиуса — остаются открытыми (Тао доказал лишь
логарифмически-усреднённый и нечётный случаи). Мы покрываем лишь коллапс-моду, и сводка названа
\code{chowla\_locked\_behind\_flip\_wall}, без слова «двигатель»; это не решение Чоулы и не Гёдель.

\paragraph{abc.}
Гипотеза abc: для взаимно простых $a+b=c$ верно
$c<K_\varepsilon\cdot\operatorname{rad}(abc)^{1+\varepsilon}$. Полиномиальный аналог — теорема
Мейсона–Стотерса~\cite{mason1984,stothers1981} — доказан и есть в \code{mathlib}.

\begin{theorem}[\code{polynomial\_abc\_shadow}]\label{thm:polynomial_abc_shadow}
\green{} Для взаимно простых многочленов $a+b+c=0$ над полем степень ограничена радикалом
произведения (либо все производные нулевые). Это \code{Polynomial.abc} из \code{mathlib}.
\end{theorem}

\emph{Идея доказательства.} Готовая теорема \code{mathlib} (вронскиан плюс взаимная простота).
Степенная граница через радикал — и есть engine-прочтение: качество тройки не может убежать в
бесконечность, оно конечно. Радикал целого числа (\code{Int.radical}, \code{Nat.radical}) реален, не
заглушка. Гейт \code{AbcConjecture} (\red{}) остаётся открытым — заявка Мочидзуки (IUT) сообществом
не принята. Нашу эпистемику мы строим лишь вокруг наивной формы $\varepsilon=0$, опровергнутой одним
свидетелем (\code{abc\_exponent\_one\_counterexample}: $1+8=9$, $\operatorname{rad}(72)=6<9$) с
неограниченной поставкой качества (\code{unboundedQualitySupplyExponentOne}); настоящий гейт
$\varepsilon>0$ этим не тронут, вся открытость живёт в $\varepsilon$-слабине.

\paragraph{Бил и Ферма–Каталан.}
Метод бесконечного спуска Ферма — буквально наш движок, и его полиномиальная форма доказана.

\begin{theorem}[\code{polynomial\_fermat\_catalan\_shadow}]\label{thm:polynomial_fermat_catalan_shadow}
\green{} При $1/p+1/q+1/r<1$ уравнение $u\,a^p+v\,b^q+w\,c^r=0$ для взаимно простых многочленов
вынуждает $a,b,c$ быть константами. Это \code{Polynomial.flt\_catalan} из \code{mathlib}.
\end{theorem}

\emph{Идея доказательства.} Готовая теорема \code{mathlib}, её доказательство — бесконечный спуск по
характеристике. Рядом стоят \code{fermatLastTheoremFour}/\code{Three} (\code{flt\_four\_is\_descent}),
доказанные тем же спуском; мы их предъявляем и выводим descent-без-двигателя через
\code{EPMI}/\code{UniversalEngine}. Гейты \code{BealConjecture} и \code{FermatCatalanConjecture}
(\red{}) — интегральный Бил и полный Ферма–Каталан — остаются открытыми (Дармон–Гранвиль доказали
конечность лишь для фиксированных показателей, через Фалтингса). Диагональные классы закрыты спуском
(\code{beal\_no\_solution\_exponent\_three}, \code{beal\_no\_solution\_exponent\_four}), а взаимная
простота несуща (\code{beal\_common\_factor\_witness}: $2^3+2^3=2^4$).

\paragraph{Лемер.}
Мера Малера — это высота, и тот же Нортокотт, что вычислял ранг у BSD, работает здесь.

\begin{theorem}[\code{mahler\_northcott\_shadow}]\label{thm:mahler_northcott_shadow}
\green{} Многочленов ограниченной степени и ограниченной меры Малера конечно много. Это
\code{finite\_mahlerMeasure\_le} из \code{mathlib}.
\end{theorem}

\emph{Идея доказательства.} Это Нортокотт: высота фундирована, бесконечного высотного спуска нет — то
самое ядро, что у BSD и универсального движка. Граница снизу — Кронекер (\code{kronecker\_boundary}):
мера Малера равна $1$ тогда и только тогда, когда корни — корни из единицы. Гейт
\code{LehmerConjecture} (\red{}) — равномерный зазор $c>1$ над единицей для нециклотомических —
остаётся открытым (наименьшее известное — число Лемера $\approx 1.17628$). Лемер при ограниченной
степени уже зелён (\code{lehmer\_at\_bounded\_degree}): для каждой степени есть $c>1$ без мер в
$(1,c)$; но $c=c(n)$ зависит от степени, а единая $c$ на все степени (\code{UniformLehmerGap},
\red{}) — и есть сама гипотеза (\code{uniformLehmerGap\_iff\_all\_degrees}).

\paragraph{Landau: простые $n^2+1$.}
Дирихле (в \code{mathlib}) даёт простые в арифметических прогрессиях, на которых стоят стороны
$6m\pm1$. Но $n^2+1$ — не прогрессия, и его бесконечность остаётся гейтом.

\begin{theorem}[\code{landauPrimes\_infinite\_of\_unbounded}]\label{thm:landauPrimes_infinite_of_unbounded}
\green{} Если простых $n^2+1$ неограниченно много, то их множество бесконечно. Кроме того
(\code{oddLandauPrime\_even\_k}), при нечётном $k$ число $k^2+1$ чётно, так что простое $k^2+1>2$
требует чётного $k$.
\end{theorem}

\emph{Идея доказательства.} Чистая структурная переупаковка (как в ветви Полиньяка): из
неограниченности следует бесконечность, без нового содержания — вся суть в гейте. Гейт
\code{Landau4thUnbounded} (\red{}) — бесконечность простых $n^2+1$ (Харди–Литтлвуд) — остаётся
открытым (Иванец: бесконечно много $n^2+1$ с не более чем двумя простыми множителями). Гейт
локализован на чётный канал (\code{landauPrime\_even\_of\_two\_lt}), а простые Ландау лежат в реальном
дирихлетовском классе $1\bmod 4$, бесконечном по \code{mathlib}
(\code{landau\_lives\_in\_dirichlet\_class}); открыта селекция внутри шеренги. Здесь, в отличие от
Била, мы построили двигательный факт (\code{landauRefutation\_carries\_engine}), хотя лежащий в основе
закон \code{LandauManifestationLaw} не декретирован, а \code{landauCause\_unknowable} — не решение
4-й проблемы и не Гёдель.

\subsection{Общая мораль}

Пять теней — один приём: где у людей есть доказанный якорь (полиномиальный abc и Ферма–Каталан,
Нортокотт и Кронекер, Дирихле, структура Лиувилля), мы предъявляем его зелёно как прочтение через
движок; где стоит открытая гипотеза — называем её честным гейтом и не притворяемся, что закрыли. Ни
одна из пяти не решена, декретов первопричины не добавлено, и один и тот же запрет узнаётся и за
горизонтом восьми масок.
